% tests/test_signal_docs.py reads the macros below, so nest braces at most one level deep in their arguments.
\documentclass[11pt]{article}

\usepackage[margin=1in]{geometry}
\usepackage{amsmath,amssymb}
\usepackage{booktabs}
\usepackage{tabularx}
\usepackage[hidelinks]{hyperref}

\newcommand{\param}[3]{#1 & \texttt{#2} & #3 \\}
\newcommand{\missingreason}[2]{\item \texttt{#1}: #2}
\newcommand{\eventaction}[1]{\texttt{#1}}
\newcommand{\signalcolumn}[2]{\texttt{#1} & #2 \\}
\newcommand{\exrebalance}[1]{#1}
\newcommand{\exparam}[2]{\texttt{#1} $=$ #2}
\newcommand{\exfiling}[6]{#1 & #2 & #3 & #4 & #5 & #6 \\}
\newcommand{\exevent}[3]{#1 & #2 & \texttt{#3} \\}
\newcommand{\exresult}[7]{#1 & #2 & #3 & #4 & #5 & #6 & #7 \\}

\newcommand{\glossaryterm}[2]{\item[#1] #2}

\title{Percentage Shares Issued\\\large Methodology}
\date{}

\begin{document}
\maketitle

\noindent This note describes how the Percentage Shares Issued Signal is measured and how
Constituents are ranked on it. It describes the method only: it holds no results, and no
parameter values beyond those of the illustrative example in Appendix~\ref{app:example}.
The code in \texttt{src/signals.py} is the source of truth;
\texttt{tests/test\_signal\_docs.py} fails when this note and the code disagree.

\section{Terms}

Capitalised terms follow \texttt{GLOSSARY.md}. The ones this note relies on, quoted from it:

\begin{description}
  \glossaryterm{Constituent}{A company in the S\&P~500 on a given date, judged by what the index held on that date and not by what it holds today.}
  \glossaryterm{Share Count}{The number of shares outstanding reported on a filing's cover page, adjusted for later stock splits.}
  \glossaryterm{Company Share Count}{The Share Count times its share factor: the Company's shares across all Share Classes, in units of its Primary Class. Signals use it in place of the Share Count.}
  \glossaryterm{Filing Date}{The date a Company's report reached the SEC.}
  \glossaryterm{Effective Date}{The first trading day after the Filing Date, from which a Share Count may be used.}
  \glossaryterm{Re-filing}{A later \texttt{ARQ} filing for a quarter the Company has already filed, often an amendment. By default the Panel keeps the first filing's Share Count and drops the Re-filing.}
  \glossaryterm{Corporate Event}{An acquisition, merger, spin-off or stock split that can change a Share Count or a price without routine issuance or buybacks.}
  \glossaryterm{Rebalance Date}{A date on which Signals are measured and Companies are ranked.}
  \glossaryterm{Lookback}{The span between the two filings a change-based Signal compares, measured between their period ends.}
  \glossaryterm{Net Shares Issued}{The change in the Company Share Count over the Lookback, counted in shares. It cannot be compared across Companies, so it is never ranked.}
  \glossaryterm{Percentage Shares Issued}{The change in the Company Share Count over the Lookback as a fraction of the earlier Company Share Count. It is a Signal, and a negative value means a net buyback.}
\end{description}

\section{Parameters}

\begin{center}
\begin{tabularx}{\linewidth}{l l X}
\toprule
Symbol & Argument & Meaning \\
\midrule
\param{$F$}{frequency}{The calendar period between Rebalance Dates, such as a quarter or a month.}
\param{$L$}{lookback\_months}{The Lookback, in months.}
\param{$\tau$}{tolerance\_days}{How far, in days, the comparison filing's period end may sit from its target.}
\param{$A$}{max\_filing\_age\_days}{The oldest, in days, the latest filing's period end may be before a Constituent counts as stale.}
\param{$x$}{exclude\_stale}{Whether stale Constituents are left unranked.}
\param{$\mathcal{B}$}{scope\_breaks}{For each Company, the dates on which the Share Classes its Share Count covers change.}
\bottomrule
\end{tabularx}
\end{center}

\section{Method}

\subsection{Rebalance Dates}

Split the trading days of the data into calendar periods of frequency $F$. A period is
\emph{complete} once the last trading day in the data falls on or after the period's last
calendar day. Each complete period contributes one Rebalance Date $t$: its last trading day.
On each $t$, the Constituents are the Companies in the index on that day.

\subsection{Filings}

For Company $i$, each filing $f$ carries a Filing Date $d_f$, an Effective Date $e_f$, a
period end $p_f$, a Share Count $s_f$ and a share factor $\phi_f$. Its Company Share Count is
\begin{equation}
  S_f = s_f \, \phi_f .
\end{equation}

\paragraph{Latest filing.} The latest filing of $i$ on $t$ is the one with the most recent
Effective Date on or before $t$:
\begin{equation}
  f_t = \operatorname*{arg\,max}_{f \,:\, e_f \le t} e_f .
\end{equation}
A filing reaches the SEC before the Rebalance Date but becomes usable only from its Effective
Date, so no Share Count is used before the market could have seen it.

\paragraph{Comparison filing.} The comparison filing is the one whose period end sits closest
to one Lookback before the latest filing's period end, among filings already effective on $t$
and within the tolerance:
\begin{align}
  p^{*} &= p_{f_t} - L \text{ months}, \\
  K_t &= \{\, g : e_g \le t,\ |p_g - p^{*}| \le \tau \,\}, \\
  g_t &= \operatorname*{arg\,min}_{g \in K_t} |p_g - p^{*}| .
\end{align}
When two filings in $K_t$ are equally close, $g_t$ is the one with the earlier Filing Date. So
when a quarter was re-filed, it keeps the first Share Count (ADR~0001).

\subsection{Net and Percentage Shares Issued}

\begin{align}
  N_{i,t} &= S_{f_t} - S_{g_t}, \\
  P_{i,t} &= \frac{N_{i,t}}{S_{g_t}} .
\end{align}
$N_{i,t}$ is counted in shares and so cannot be compared across Companies; it is never ranked.
$P_{i,t}$ is the Signal. A negative value is a net buyback.

The filing age is the number of days from the latest period end to the Rebalance Date,
$a_{i,t} = t - p_{f_t}$.

\subsection{Constituents with no value}

A Constituent with no value stays in the table, holds no rank, and records the first reason
that applies, in this priority order:

\begin{enumerate}
  \missingreason{no\_share\_count}{no filing is effective on $t$, or its Company Share Count is missing (ADR~0003).}
  \missingreason{stale}{$x$ is set and $a_{i,t} > A$.}
  \missingreason{no\_comparison}{$K_t$ is empty, or $S_{g_t}$ is missing or $\le 0$.}
  \missingreason{scope\_break}{a date $b \in \mathcal{B}_i$ falls in $d_{g_t} < b \le d_{f_t}$, so the change reflects which Share Classes are counted rather than issuance or buybacks (ADR~0002).}
\end{enumerate}

\subsection{Rank}

Let $V_t$ be the Constituents with a value on $t$, $n = |V_t|$, and $r_{i,t}$ the rank of
$P_{i,t}$ within $V_t$ in ascending order, tied values sharing the average of their ranks.
Then
\begin{equation}
  R_{i,t} =
  \begin{cases}
    \dfrac{r_{i,t} - 1}{n - 1} & n \ge 2, \\[1ex]
    \tfrac{1}{2} & n = 1 .
  \end{cases}
\end{equation}
$R = 0$ is the lowest $P$, the biggest net buyback when $P < 0$, and $R = 1$ the highest. Each Rebalance Date is
ranked on its own.

\subsection{Corporate Event flag}

A Constituent is flagged when a Corporate Event of its own falls between the Filing Dates of
its comparison and latest filings: $d_{g_t} < \text{event date} \le d_{f_t}$. The window runs
between Filing Dates, not the period ends that measure the Lookback. The flag marks changes that may
come from a deal rather than routine issuance or buybacks; it does not change the value or the
rank. The actions that count are
\eventaction{spinoff}, \eventaction{spinoffdividend}, \eventaction{spunofffrom},
\eventaction{mergerfrom}, \eventaction{mergerto} and every action starting with
\texttt{acquisition}, written \eventaction{acquisition*}. Stock splits are not flagged: the
Share Count is already adjusted for them. A Constituent with an empty $K_t$
is never flagged.

\section{Output}

One row per Constituent per Rebalance Date, with these columns in this order:

\begin{center}
\begin{tabularx}{\linewidth}{l X}
\toprule
Column & Holds \\
\midrule
\signalcolumn{date}{The Rebalance Date $t$.}
\signalcolumn{permaticker}{The Company's permanent identifier.}
\signalcolumn{ticker}{The Company's Primary Class ticker.}
\signalcolumn{net\_shares\_issued}{$N_{i,t}$, empty when the Constituent has no value.}
\signalcolumn{pct\_shares\_issued}{$P_{i,t}$, empty when the Constituent has no value.}
\signalcolumn{pct\_shares\_issued\_rank}{$R_{i,t}$, empty when the Constituent has no value.}
\signalcolumn{filing\_age\_days}{$a_{i,t}$, empty when no filing is effective.}
\signalcolumn{missing\_reason}{The reason the Constituent has no value, or empty.}
\signalcolumn{corporate\_event\_in\_lookback}{The Corporate Event flag.}
\bottomrule
\end{tabularx}
\end{center}

\section{Data caveats}

These follow from the data, whatever the parameters.

\begin{itemize}
  \item The earliest Filing Date in the data is 2016-09-30. A Constituent has a comparison
    filing only once its latest period end is at least $L$ months, less $\tau$, after the
    earliest period end on record, so the first Rebalance Dates rank few or no Constituents.
  \item The Panel carries the filings in force on the trading days it covers. A filing that was
    replaced before a Company's first trading day in the Panel is not available as a
    comparison filing, so a Company whose listing starts late can lack one for its first
    Rebalance Dates.
\end{itemize}

\clearpage
\appendix
\section{Illustrative example}
\label{app:example}

Synthetic Companies A, B and C on one Rebalance Date, \exrebalance{Y2-12-31}. Dates fall in
years Y1 and Y2, neither a leap year, and each Effective Date is taken as the day after the
Filing Date. The example uses \exparam{lookback\_months}{12}, \exparam{tolerance\_days}{40},
\exparam{max\_filing\_age\_days}{200} and \exparam{exclude\_stale}{true}, with no scope breaks.

\begin{itemize}
  \item A issues shares and reports them across two Share Classes (share factor 2). It also
    makes an acquisition in the Lookback.
  \item B buys back shares. Its comparison quarter was re-filed with a different Share Count,
    which is ignored.
  \item C's latest filing is 275 days old, so C is stale. Its stock split is not a flagged
    Corporate Event.
\end{itemize}

\paragraph{Filings.}
\begin{center}
\begin{tabular}{l l l l r r}
\toprule
Company & Period end & Filing Date & Effective Date & Share Count & Share factor \\
\midrule
\exfiling{A}{Y1-09-30}{Y1-10-30}{Y1-10-31}{500}{2}
\exfiling{A}{Y2-09-30}{Y2-10-30}{Y2-10-31}{550}{2}
\exfiling{B}{Y1-09-30}{Y1-11-05}{Y1-11-06}{2000}{1}
\exfiling{B}{Y1-09-30}{Y1-12-14}{Y1-12-15}{2100}{1}
\exfiling{B}{Y2-09-30}{Y2-11-04}{Y2-11-05}{1900}{1}
\exfiling{C}{Y1-03-31}{Y1-05-01}{Y1-05-02}{3000}{1}
\exfiling{C}{Y2-03-31}{Y2-05-01}{Y2-05-02}{3300}{1}
\bottomrule
\end{tabular}
\end{center}

\paragraph{Corporate Events.}
\begin{center}
\begin{tabular}{l l l}
\toprule
Company & Date & Action \\
\midrule
\exevent{A}{Y2-03-15}{acquisitionby}
\exevent{C}{Y2-02-01}{split}
\bottomrule
\end{tabular}
\end{center}

\paragraph{Result.} A: $N = 1100 - 1000 = 100$, $P = 100 / 1000 = 0.10$.
B: $N = 1900 - 2000 = -100$, $P = -100 / 2000 = -0.05$, using the first filing of Y1-09-30.
C is stale, so A and B are the $n = 2$ valued Constituents: B ranks $0$ and A ranks $1$.

\begin{center}
\begin{tabular}{l r r r r l l}
\toprule
Company & $N$ & $P$ & $R$ & $a$ & Missing reason & Event \\
\midrule
\exresult{A}{100}{0.10}{1}{92}{---}{yes}
\exresult{B}{-100}{-0.05}{0}{92}{---}{no}
\exresult{C}{---}{---}{---}{275}{stale}{no}
\bottomrule
\end{tabular}
\end{center}

\end{document}
